\documentclass[10pt,aspectratio=169]{beamer}
% Preámbulo modular para presentaciones Beamer (Modelación Estadística)
% Diseño y paleta institucional del Tecnológico de Monterrey

\documentclass[aspectratio=169,xcolor={dvipsnames,table}]{beamer}

\usepackage[utf8]{inputenc}
\usepackage[spanish,mexico]{babel}
\usepackage{amsmath,amssymb,amsfonts,mathtools}
\usepackage{graphicx}
\usepackage{listings}
\usepackage{booktabs}
\usepackage{tikz}
\usetikzlibrary{matrix,arrows,decorations.pathmorphing,shapes.geometric,positioning}

% Paleta Institucional Tecnológico de Monterrey
\definecolor{TecAzul}{HTML}{0039A6}       % Azul TEC: color primario institucional
\definecolor{TecAzulOscuro}{HTML}{1A2E51} % Azul marino oscuro
\definecolor{TecRojo}{HTML}{EC2661}       % Rojo de acento (paleta miTec)
\definecolor{TecGris}{HTML}{646464}       % Gris neutro para comentarios y subtítulos
\definecolor{TecFondoBloque}{HTML}{F4F6F9}% Fondo suave para bloques

% Configuración de Tema Beamer
\usetheme{Boadilla}
\usecolortheme[named=TecAzulOscuro]{structure}
\setbeamercolor{alerted text}{fg=TecRojo}
\setbeamercolor{palette primary}{bg=TecAzulOscuro,fg=white}
\setbeamercolor{palette secondary}{bg=TecAzul,fg=white}
\setbeamercolor{palette tertiary}{bg=TecAzulOscuro!80!black,fg=white}
\setbeamercolor{title}{fg=TecAzulOscuro,bg=TecFondoBloque}
\setbeamercolor{frametitle}{fg=TecAzulOscuro,bg=TecFondoBloque}

% Bloques personalizados elegantes
\setbeamercolor{block title}{bg=TecAzulOscuro,fg=white}
\setbeamercolor{block body}{bg=TecFondoBloque,fg=black}
\setbeamercolor{block title alerted}{bg=TecRojo,fg=white}
\setbeamercolor{block body alerted}{bg=TecRojo!8,fg=black}
\setbeamercolor{block title example}{bg=TecAzul,fg=white}
\setbeamercolor{block body example}{bg=TecAzul!8,fg=black}

% Redondear bordes y añadir sombra sutil
\setbeamertemplate{blocks}[rounded][shadow=true]
\setbeamertemplate{navigation symbols}{}

% Pie de página limpio con número de diapositiva
\setbeamertemplate{footline}{
  \leavevmode%
  \hbox{%
  \begin{beamercolorbox}[wd=.7\paperwidth,ht=2.25ex,dp=1ex,left]{author in head/foot}%
    \hspace*{2ex}\insertshortauthor~~(\insertshortinstitute) \hfill \insertshorttitle
  \end{beamercolorbox}%
  \begin{beamercolorbox}[wd=.3\paperwidth,ht=2.25ex,dp=1ex,right]{date in head/foot}%
    \insertshortdate{}\hspace*{2em}
    \insertframenumber{} / \inserttotalframenumber\hspace*{2ex}
  \end{beamercolorbox}}%
  \vskip0pt%
}

% Configuración de Listings de Python para visualización pedagógica de código
\lstset{
    language=Python,
    basicstyle=\ttfamily\footnotesize,
    keywordstyle=\color{TecAzulOscuro}\bfseries,
    stringstyle=\color{TecRojo},
    commentstyle=\color{TecGris}\itshape,
    showstringspaces=false,
    breaklines=true,
    frame=lines,
    rulecolor=\color{TecAzulOscuro!30},
    backgroundcolor=\color{TecFondoBloque},
    aboveskip=0.5em,
    belowskip=0.5em,
    literate={á}{{\'a}}1 {é}{{\'e}}1 {í}{{\'i}}1 {ó}{{\'o}}1 {ú}{{\'u}}1
             {Á}{{\'A}}1 {É}{{\'E}}1 {Í}{{\'I}}1 {Ó}{{\'O}}1 {Ú}{{\'U}}1
             {ñ}{{\~n}}1 {Ñ}{{\~N}}1 {¿}{{?`}}1 {¡}{{!`}}1
}

% Metadatos institucionales por defecto
\title[Modelación Estadística]{Modelación Estadística para Ciencia de Datos e Ingeniería}
\author[J. Castillo Colmenares]{Juliho David Castillo Colmenares}
\institute[Tec de Monterrey]{Tecnológico de Monterrey \\ Escuela de Ingeniería y Ciencias}
\date{\today}

% Comandos y macros estadísticas compartidas para las presentaciones Beamer (Metropolis)

% Conjuntos numéricos básicos
\newcommand{\R}{\mathbb{R}}
\newcommand{\N}{\mathbb{N}}
\newcommand{\Q}{\mathbb{Q}}
\newcommand{\Z}{\mathbb{Z}}
\newcommand{\set}[1]{\left\{ #1 \right\}}
\newcommand{\abs}[1]{\left|#1\right|}
\newcommand{\norm}[1]{\left\| #1 \right\|}

% Comandos estadísticos y probabilísticos del curso
\newcommand{\Var}{\operatorname{Var}}
\newcommand{\cov}{\operatorname{Cov}}
\newcommand{\corr}{\rho}
\newcommand{\comb}[2]{\begin{pmatrix} #1 \\ #2 \end{pmatrix}}
\newcommand{\s}{\sigma}
\newcommand{\particion}{\mathfrak{P}}
\newcommand{\card}[1]{\#\left( #1 \right)}
\newcommand{\seq}[3]{\set{#1}_{#2}^{#3}}
\newcommand{\E}{\mathbb{E}}
\newcommand{\Prob}{\mathbb{P}}

% Entornos pedagógicos y cajas en español académico natural
\newenvironment{discusion}{%
  \begin{alertblock}{Pregunta de Discusión}%
}{%
  \end{alertblock}%
}

\newenvironment{reflexion}{%
  \begin{alertblock}{Reflexión para el Aula}%
}{%
  \end{alertblock}%
}

\newenvironment{paradoja}{%
  \begin{alertblock}{Paradoja de Exploración}%
}{%
  \end{alertblock}%
}

\newenvironment{motivacion}{%
  \begin{exampleblock}{Motivación: Ciencia de Datos en la Práctica}%
}{%
  \end{exampleblock}%
}

\newenvironment{retoclase}{%
  \begin{block}{Reto Rápido en Equipo}%
}{%
  \end{block}%
}


\title[Técnicas de conteo]{Sección 02.03: Técnicas de conteo}
\subtitle{Permutaciones, combinaciones y probabilidad clásica}
\author[J. Castillo Colmenares]{Juliho Castillo Colmenares}
\institute{Tecnológico de Monterrey}
\date{\vspace{-1.2cm}}

\begin{document}

% 01 — Identidad
\begin{frame}[plain]
  \vspace{-0.8cm}
  \titlepage
\end{frame}

% 02 — Hoja de ruta
\begin{frame}{Hoja de ruta --- Capítulo 02: Teoría de probabilidad}
  \scriptsize
  \begin{itemize}\itemsep=0.04em
    \item \textbf{02.01} Teoría de conjuntos y probabilidad.
    \item \textbf{02.02} Fundamentos y axiomas de probabilidad.
    \item \textbf{02.03} Técnicas de conteo \emph{(hoy)}.
    \item \textbf{02.04} Probabilidad condicional y regla de Bayes.
    \item \textbf{02.05} Teorema de Bayes.
    \item \textbf{02.06} Muestreo aleatorio.
  \end{itemize}
  \vspace{-0.05cm}
  \begin{block}{Objetivo didáctico}
    Contar espacios muestrales finitos usando el principio de multiplicación,
    permutaciones y combinaciones, y convertir esos conteos en probabilidades
    clásicas cuando los resultados son equiprobables.
  \end{block}
\end{frame}

% 03 — Motivación
\begin{frame}{De enumerar resultados a contar sistemáticamente}
  \small
  \begin{columns}[T]
    \begin{column}{0.48\textwidth}
      \begin{block}{El reto}
        Un espacio muestral puede crecer demasiado rápido para listar cada
        resultado. El conteo permite describirlo sin enumerarlo.
      \end{block}
    \end{column}
    \begin{column}{0.48\textwidth}
      \begin{alertblock}{Probabilidad clásica}
        Si todos los resultados son igualmente probables,
        \[
          P(A)=\frac{\lvert A\rvert}{\lvert\Omega\rvert}.
        \]
        Las técnicas de conteo calculan ambos cardinales.
      \end{alertblock}
    \end{column}
  \end{columns}
\end{frame}

% 04 — Principio fundamental
\begin{frame}{Principio fundamental de conteo}
  \small
  \begin{block}{Regla del producto}
    Si una primera operación tiene $n_1$ opciones, una segunda tiene $n_2$
    opciones para cada elección anterior, y así sucesivamente, entonces
    \[
      N=n_1n_2\cdots n_k.
    \]
  \end{block}
  \pause
  \begin{exampleblock}{Interpretación}
    Cada resultado final es una ruta: se multiplica porque cada etapa se
    combina con todas las opciones de las etapas posteriores.
  \end{exampleblock}
\end{frame}

% 05 — Modelado del conteo
\begin{frame}{Cómo modelar un conteo}
  \footnotesize
  \begin{enumerate}\itemsep=0.12cm
    \item Describir con precisión el resultado que se cuenta.
    \item Separar el proceso en etapas sucesivas o elecciones.
    \item Preguntar si el orden importa y si se permite repetir elementos.
    \item Elegir multiplicación, permutación o combinación.
  \end{enumerate}
  \vspace{0.1cm}
  \begin{alertblock}{Regla de control}
    Dos resultados que cambian al intercambiar posiciones requieren un modelo
    ordenado; si no cambian, se cuenta una selección sin orden.
  \end{alertblock}
\end{frame}

% 06 — Permutaciones
\begin{frame}{Permutaciones: cuando el orden importa}
  \small
  \begin{block}{Definición}
    El número de arreglos ordenados de $r$ objetos distintos tomados de $n$,
    sin repetición, es
    \[
      P(n,r)=\frac{n!}{(n-r)!}=n(n-1)\cdots(n-r+1).
    \]
  \end{block}
  \begin{exampleblock}{Caso completo}
    Si $r=n$, se obtiene $P(n,n)=n!$. El orden distingue cada arreglo.
  \end{exampleblock}
\end{frame}

% 07 — Combinaciones
\begin{frame}{Combinaciones: cuando el orden no importa}
  \small
  \begin{block}{Definición}
    El número de subconjuntos de tamaño $r$ tomados de $n$ objetos es
    \[
      \binom{n}{r}=\frac{P(n,r)}{r!}=\frac{n!}{r!(n-r)!}.
    \]
  \end{block}
  \begin{exampleblock}{Idea clave}
    Cada selección sin orden produce $r!$ arreglos ordenados; dividir por
    $r!$ elimina esa multiplicidad.
  \end{exampleblock}
\end{frame}

% 08 — Propiedades y decisión
\begin{frame}{Propiedades y decisión entre modelos}
  \footnotesize
  \begin{columns}[T]
    \begin{column}{0.48\textwidth}
      \begin{block}{Coeficiente binomial}
        \[
          \binom{n}{r}=\binom{n}{n-r},\qquad
          \binom{n}{0}=\binom{n}{n}=1.
        \]
      \end{block}
    \end{column}
    \begin{column}{0.48\textwidth}
      \begin{alertblock}{Pregunta diagnóstica}
        \begin{itemize}\itemsep=0.05cm
          \item ¿Cambiar el orden cambia el resultado? $\to P(n,r)$.
          \item ¿Solo importa el conjunto elegido? $\to \binom{n}{r}$.
          \item ¿Hay etapas sucesivas? $\to n_1\cdots n_k$.
        \end{itemize}
      \end{alertblock}
    \end{column}
  \end{columns}
\end{frame}

% 09 — Puente computacional 1/4
\begin{frame}{Puente computacional 1/4: aritmética exacta}
  \small
  \begin{block}{Verificación reproducible}
    Las expresiones del tema pueden verificarse con enteros, sin redondear:
    \[
      8!=40320,\qquad 5!=120,\qquad 10!=3628800.
    \]
  \end{block}
  \begin{exampleblock}{Principio de control}
    Primero se calcula el factorial; después se simplifica la razón. Así se
    distingue un error de conteo de un error de aproximación numérica.
  \end{exampleblock}
\end{frame}

% 10 — Puente computacional 2/4
\begin{frame}{Puente computacional 2/4: permutaciones y combinaciones}
  \footnotesize
  \begin{columns}[T]
    \begin{column}{0.48\textwidth}
      \begin{block}{Permutación}
        \[
          P(8,3)=8\cdot7\cdot6=336.
        \]
        El orden de oro, plata y bronce produce resultados distintos.
      \end{block}
    \end{column}
    \begin{column}{0.48\textwidth}
      \begin{block}{Combinación}
        \[
          \binom{10}{4}=\frac{10!}{4!6!}=210.
        \]
        Intercambiar integrantes no crea un comité nuevo.
      \end{block}
    \end{column}
  \end{columns}
\end{frame}

% 11 — Puente computacional 3/4
\begin{frame}{Puente computacional 3/4: espacio muestral}
  \small
  \begin{block}{Manos de cinco cartas}
    Para una mano de $5$ cartas de una baraja de $52$, el espacio muestral es
    \[
      \lvert\Omega\rvert=\binom{52}{5}=2{,}598{,}960.
    \]
    El orden de recepción no cambia la mano final.
  \end{block}
  \pause
  \begin{alertblock}{Casos favorables}
    Un palo tiene $\binom{13}{5}$ manos y hay cuatro palos, de modo que
    $\lvert A\rvert=4\binom{13}{5}=5{,}148$.
  \end{alertblock}
\end{frame}

% 12 — Puente computacional 4/4
\begin{frame}{Puente computacional 4/4: comprobaciones numéricas}
  \scriptsize
  \begin{center}
    \begin{tabular}{lrr}
      \hline
      Conteo & Resultado exacto & Interpretación \\
      \hline
      $4\cdot6\cdot3$ & $72$ & menús \\
      $P(8,3)$ & $336$ & medallas ordenadas \\
      $\binom{10}{4}$ & $210$ & comité sin cargos \\
      $4\binom{13}{5}$ & $5{,}148$ & manos del mismo palo \\
      \hline
    \end{tabular}
  \end{center}
  \vspace{0.1cm}
  Cada valor se obtiene con las fórmulas de la sección y puede reproducirse
  con aritmética entera independiente del idioma del mazo.
\end{frame}

% 13–14 — Ejemplo 1
\begin{frame}{Ejemplo 1/4 --- enunciado: menú de restaurante}
  \small
  Un restaurante ofrece $4$ entradas, $6$ platos fuertes y $3$ postres.
  Si se elige exactamente una opción de cada categoría, ¿cuántos menús
  distintos pueden formarse?
  \vfill
  \begin{alertblock}{Modelo}
    Tres etapas sucesivas; todas las opciones se combinan con las posteriores.
  \end{alertblock}
\end{frame}

\begin{frame}{Ejemplo 1/4 --- resolución}
  \small
  \begin{block}{Aplicación del principio de multiplicación}
    \[
      N=4\cdot6\cdot3=72.
    \]
  \end{block}
  \begin{exampleblock}{Interpretación}
    El conteo no enumera los menús: construye cada menú como una ruta de tres
    decisiones y cuenta todas las rutas posibles.
  \end{exampleblock}
\end{frame}

% 15–16 — Ejemplo 2
\begin{frame}{Ejemplo 2/4 --- enunciado: medallas}
  \small
  En una carrera con $8$ corredores, ¿de cuántas formas pueden otorgarse las
  medallas de oro, plata y bronce, sin empates?
  \vfill
  \begin{alertblock}{Diagnóstico}
    Una asignación de medallas es ordenada: cambiar el puesto de dos
    corredores produce una asignación diferente.
  \end{alertblock}
\end{frame}

\begin{frame}{Ejemplo 2/4 --- resolución}
  \small
  \begin{block}{Permutación de $8$ tomados de $3$}
    \[
      P(8,3)=\frac{8!}{(8-3)!}=8\cdot7\cdot6=336.
    \]
  \end{block}
  \begin{exampleblock}{Decisión}
    Se usa $P(n,r)$ porque oro, plata y bronce no son posiciones
    intercambiables.
  \end{exampleblock}
\end{frame}

% 17–18 — Ejemplo 3
\begin{frame}{Ejemplo 3/4 --- enunciado: comité}
  \small
  Se forma un comité de $4$ personas a partir de $10$ candidatos. Todos los
  puestos son equivalentes. ¿Cuántos comités distintos pueden formarse?
  \vfill
  \begin{alertblock}{Diagnóstico}
    El comité es un subconjunto: el orden en que se eligen sus integrantes no
    cambia el resultado.
  \end{alertblock}
\end{frame}

\begin{frame}{Ejemplo 3/4 --- resolución}
  \small
  \begin{block}{Combinación de $10$ tomados de $4$}
    \[
      \binom{10}{4}=\frac{10!}{4!6!}=210.
    \]
  \end{block}
  \begin{exampleblock}{Interpretación}
    Las $4!$ maneras de ordenar cada comité representan el mismo comité y se
    eliminan mediante la fórmula de combinaciones.
  \end{exampleblock}
\end{frame}

% 19–20 — Ejemplo 4
\begin{frame}{Ejemplo 4/4 --- enunciado: mano del mismo palo}
  \small
  De una baraja estándar de $52$ cartas se extrae una mano de $5$. ¿Cuál es
  la probabilidad de que las cinco cartas sean del mismo palo?
  \vfill
  \begin{alertblock}{Modelo}
    El espacio muestral y los casos favorables son combinaciones, porque el
    orden de recepción no importa.
  \end{alertblock}
\end{frame}

\begin{frame}{Ejemplo 4/4 --- resolución}
  \footnotesize
  \begin{block}{Probabilidad clásica}
    \[
      P(A)=\frac{4\binom{13}{5}}{\binom{52}{5}}
      =\frac{5{,}148}{2{,}598{,}960}\approx0.00198.
    \]
  \end{block}
  \begin{exampleblock}{Lectura del resultado}
    Aproximadamente una de cada $505$ manos tiene las cinco cartas del mismo
    palo; el denominador cuenta todas las manos y el numerador solo las
    favorables.
  \end{exampleblock}
\end{frame}

% 21 — Síntesis
\begin{frame}{Síntesis de la sección}
  \small
  \begin{itemize}\itemsep=0.12cm
    \item La regla del producto cuenta etapas sucesivas.
    \item Las permutaciones modelan resultados ordenados.
    \item Las combinaciones modelan selecciones sin orden.
    \item La probabilidad clásica es el cociente entre casos favorables y
      resultados equiprobables.
  \end{itemize}
\end{frame}

% 22 — Transición
\begin{frame}{Transición: del conteo a la probabilidad condicional}
  \small
  \begin{block}{Lo que ya dominamos}
    Podemos construir espacios muestrales finitos y medir eventos cuando todos
    los resultados tienen la misma probabilidad.
  \end{block}
  \begin{alertblock}{Siguiente sección --- 02.04}
    La probabilidad condicional restringe el espacio muestral a la información
    disponible y prepara la regla de Bayes.
  \end{alertblock}
\end{frame}

\end{document}
